\documentclass[a4paper,12pt]{article}
\usepackage[utf8]{inputenc}
\usepackage[english, russian]{babel}
\usepackage{amsmath, amssymb, amsthm}
\usepackage{enumitem}%оформление списков
\usepackage{varwidth}
\usepackage{graphicx}
\usepackage{float}
\usepackage{indentfirst}
\usepackage{url}
\usepackage{seqsplit}
\usepackage{pdflscape}
\usepackage{tabularx}
\usepackage{xcolor}
\usepackage{listings}
\usepackage{mathtools}
\usepackage{hyperref}
\usepackage[backend=biber, style=numeric, sorting=none]{biblatex}

\addbibresource{my_libs.bib}


\usepackage[a4paper, top=20mm, left=30mm, right=20mm, bottom=20mm]{geometry}
\newcommand{\dim}{\mathrm{dim}}
\newcommand{\lin}{\mathrm{lin}}
\newcommand{\diag}{\mathrm{diag}}
\newcommand{\R}{\mathbb{R}}
\newcommand{\Rn}{\mathbb{R^n}}

\lstset{
    language=Python,
    basicstyle=\ttfamily\small,
    keywordstyle=\color{blue},
    commentstyle=\color{green!50!black},
    stringstyle=\color{red},
    numbers=left,
    numberstyle=\tiny,
    stepnumber=1,
    showstringspaces=false,
    breaklines=true,
    frame=single,
    captionpos=b
}

\newtheoremstyle{indented}% отступ
{3pt}% над
{3pt}% под
{\itshape}% шрифт тела, переключатель
{\parindent}% отступ
{\bfseries}% шрифт заголовка
{.}% точка после заголовка
{.5em}% отступ от заголовка
{}% head spec (empty = ‘normal’)

\theoremstyle{indented}
\newtheorem{theorem}{Теорема}
\newtheorem*{theorem*}{Теорема}
\newtheorem{definition}{Определение}
\newtheorem*{definition*}{Определение}
\newtheorem{following}{Следствие}
\newtheorem*{following*}{Следствие}
\newtheorem{statement}{Утверждение}
\newtheorem*{statement*}{Утверждение}

\renewcommand{\thesection}{\arabic{section}.}
\renewcommand{\thesubsection}{\arabic{section}.\arabic{subsection}.}
\renewcommand{\thesubsubsection}{\arabic{section}.\arabic{subsection}.\arabic{subsubsection}.}

\setcounter{section}{0}
\setcounter{subsection}{0}
\setcounter{equation}{0}

\begin{document}
\pagenumbering{arabic}
\setcounter{page}{1}
\pagestyle{empty}
\begin{center}
\textbf{Министерство науки и высшего образования Российской Федерации}
\\
\vspace{0.5cm}
\textbf{Федеральное государственное бюджетное образовательное\\
учреждение высшего образования} \\
\textbf{«Ярославский государственный университет им. П.Г. Демидова»}

\vspace{0.5cm}

Кафедра математического анализа
\vfill

\begin{flushright}
  Сдано на кафедру\\
  « \underline{\hspace{0.8cm}} » \underline{\hspace{3cm}} 2026 г.\\
  Заведующий кафедрой \\
  \underline{д. ф.-м. н., доцент }\\
  \underline{\hspace{3.4cm}} Невский М.В.\\
\end{flushright}
\vfill

Выпускная квалификационная работа магистра \\

\vspace{0.5cm}

\textbf{Задача о минимальной норме интерполяционного проектора} \\
(Направление подготовки магистров 01.03.02 Прикладная математика и информатика)
\vfill


\bigskip


\end{center}

\begin{flushright}
Научный руководитель\\
\underline{к. ф.-м. н., доцент }\\
\underline{\hspace{3.4cm}} Ухалов А.Ю.\\
« \underline{\hspace{0.8cm}} » \underline{\hspace{3cm}} 2026 г.\\
\vspace{1.5cm}
Студент группы \underline{ПМИ-21МО}\\
\underline{\hspace{3.5cm}} Зайков И.В. \\
« \underline{\hspace{0.8cm}} » \underline{\hspace{3cm}} 2026 г.\\
\end{flushright}

\vspace{0.5cm}
\begin{center}
Ярославль 2026 г.
\end{center}
\newpage

\section*{Реферат}
\addcontentsline{section}{Реферат}

\newpage

\pagestyle{plain}
\tableofcontents

\newpage

% ============================================================================
\section*{Введение}
\addcontentsline{toc}{section}{Введение}

В выпускной квалификационной работе рассматривается задача о минимальной норме интерполяционного проектора. Изучается интерполяция непрерывной функции $f(x)$, определённой на кубе $Q_n = [0,1]^n$, с помощью полиномов от $n$ переменных степени не выше $1$. Для построения интерполяционного полинома необходимо задать $n+1$ узел — вершины невырожденного симплекса $S \subset Q_n$. Интерполяционный проектор $P$ ставит в соответствие функции $f$ полином $P f$, совпадающий с $f$ в узлах. Норма оператора $P$ характеризует устойчивость интерполяции к погрешностям входных данных: чем меньше норма, тем ближе интерполяционный полином к наилучшему приближению.

Минимальная норма интерполяционного проектора $\theta_n$ определяется как наименьшее возможное значение $\|P\|$ при условии, что узлы интерполяции принадлежат кубу $Q_n$ и образуют невырожденный симплекс. Задача нахождения $\theta_n$ является классической в теории приближений и имеет важное практическое значение, поскольку позволяет выбрать оптимальное расположение узлов интерполяции, обеспечивающее наибольшую устойчивость.

В настоящей работе для получения верхней оценки $\theta_n$ используется подход, основанный на рассмотрении симплексов максимального объёма, вписанных в гиперкуб $Q_n$. Такие симплексы могут быть построены из максимальных $(-1/1)$-определителей порядка $n+1$, то есть из матриц Адамара. Норма проектора для фиксированного симплекса вычисляется по формуле
\[
\|P\| = \max_{x \in \operatorname{ver}(Q_n)} \sum_{j=1}^{n+1} |\lambda_j(x)|,
\]
где $\lambda_j$ — базисные многочлены Лагранжа (барицентрические координаты). Для нахождения максимума требуется перебрать все $2^n$ вершин гиперкуба.

Ручной перебор для размерности $n = 31$ ($2^{31} \approx 2.1$ миллиарда вершин) невозможен, поэтому была разработана программа на языке C++ с использованием библиотеки Eigen для линейно-алгебраических операций. Вычисления распараллелены на $16$ потоков, что позволило завершить полный перебор за $290$ секунд (около $5$ минут).

Для верификации программной реализации был проведён ряд проверок:
\begin{itemize}
    \item проверка интерполяционных свойств $\lambda_j(x^{(k)}) = \delta_{jk}$ с точностью $10^{-10}$;
    \item проверка определителей матриц Адамара на соответствие теоретическому значению $n^{n/2}$;
    \item верификация на малых размерностях ($n = 3, 5, 7$) путём сравнения с известными из литературы результатами;
    \item повторная проверка найденного максимума после завершения перебора.
\end{itemize}

В результате работы получена новая верхняя оценка минимальной нормы интерполяционного проектора для размерности $n = 31$:
\[
\theta_{31} \le 5.000000.
\]
Теоретическая верхняя граница для данной размерности составляет $\sqrt{32} \approx 5.656854249$, что подтверждает корректность разработанного алгоритма. Полученный результат является новым и уточняет ранее известные теоретические оценки.

\newpage

% ============================================================================
\section{Используемые обозначения}

Везде далее предполагается, что $n$ — натуральное число ($n \in \mathbb{N}$). Обозначим через $Q_n$ единичный куб в $\mathbb{R}^n$:
\[
Q_n = [0,1]^n = \{(x_1, \dots, x_n) : 0 \le x_i \le 1,\ i = 1,\dots,n\}.
\]

Через $C(Q_n)$ обозначим пространство непрерывных функций $f : Q_n \to \mathbb{R}$ с равномерной нормой
\[
\|f\|_{C(Q_n)} = \max_{x \in Q_n} |f(x)|.
\]

Под $\Pi_1(\mathbb{R}^n)$ понимается совокупность многочленов от $n$ переменных степени $\le 1$, т.е. линейная оболочка
\[
\Pi_1(\mathbb{R}^n) = \operatorname{span}\{1, x_1, x_2, \dots, x_n\}.
\]

Симплекс $S \subset \mathbb{R}^n$ с вершинами $x^{(1)}, \dots, x^{(n+1)}$ обозначается
\[
S = \operatorname{conv}\{x^{(1)}, \dots, x^{(n+1)}\}.
\]

% ============================================================================
\section{Интерполяционный проектор и его норма}

\subsection{Постановка задачи интерполяции}

Пусть $Q_n = [0,1]^n$ — единичный куб в $\mathbb{R}^n$, а $C(Q_n)$ — пространство непрерывных функций $f: Q_n \to \mathbb{R}$ с равномерной нормой
\[
\|f\|_{C(Q_n)} = \max_{x \in Q_n} |f(x)|.
\]

% todo: сформулировать текст
В \cite{nevsky_ukhalov_sel_problems} описывается задача интерполяции....

Рассмотрим интерполяцию функций из $C(Q_n)$ полиномами от $n$ переменных степени не выше $1$. Пространство таких полиномов обозначим через $\Pi_1(\mathbb{R}^n)$:
\[
\Pi_1(\mathbb{R}^n) = \operatorname{span}\{1, x_1, x_2, \dots, x_n\}.
\]

Для построения интерполяционного полинома необходимо задать $n+1$ узел
\[
x^{(1)}, x^{(2)}, \dots, x^{(n+1)} \in Q_n,
\]
образующих невырожденный симплекс $S$. Требуется найти полином $p \in \Pi_1(\mathbb{R}^n)$ такой, что
\[
p(x^{(j)}) = f(x^{(j)}), \quad j = 1, \dots, n+1.
\]

\subsection{Базисные многочлены Лагранжа и барицентрические координаты}

% todo: сформулировать
В \cite{nevsky_geo} показано, как построить....

Обозначим вершины симплекса $S$ через
\[
x^{(j)} = \left(x_1^{(j)}, x_2^{(j)}, \dots, x_n^{(j)}\right), \quad 1 \le j \le n+1.
\]

Рассмотрим матрицу $A$ размера $(n+1) \times (n+1)$, в которой \textbf{строки} соответствуют вершинам симплекса, дополненным единицей в последнем столбце:
\[
A = \begin{pmatrix}
x_1^{(1)} & x_2^{(1)} & \dots & x_n^{(1)} & 1 \\
x_1^{(2)} & x_2^{(2)} & \dots & x_n^{(2)} & 1 \\
\vdots & \vdots & & \vdots & \vdots \\
x_1^{(n+1)} & x_2^{(n+1)} & \dots & x_n^{(n+1)} & 1
\end{pmatrix}.
\]

Столбец единиц введён для того, чтобы учесть условие нормировки барицентрических координат $\sum_{j=1}^{n+1} \lambda_j(x) = 1$ в виде линейного уравнения. Благодаря этому переход к однородным координатам позволяет записать систему для нахождения барицентрических координат в матричной форме.

Пусть $\Delta = \det(A)$. Поскольку симплекс $S$ невырожден, $\Delta \neq 0$. Обозначим через $\Delta_j(x)$ определитель, который получается из $\Delta$ заменой $j$-й строки на строку $(x_1, x_2, \dots, x_n, 1)$. Тогда функции
\[
\lambda_j(x) = \frac{\Delta_j(x)}{\Delta}, \quad j = 1, \dots, n+1,
\]
называются \textbf{барицентрическими координатами} точки $x$ относительно симплекса $S$. Они обладают следующими свойствами:
\begin{enumerate}
    \item $\lambda_j(x) \in \Pi_1(\mathbb{R}^n)$, то есть являются линейными функциями;
    \item $\lambda_j(x^{(k)}) = \delta_{jk}$, где $\delta_{jk}$ — символ Кронекера;
    \item $\sum_{j=1}^{n+1} \lambda_j(x) = 1$ для всех $x \in \mathbb{R}^n$.
\end{enumerate}

Таким образом, барицентрические координаты $\lambda_j(x)$ совпадают с \textbf{базисными многочленами Лагранжа}, соответствующими симплексу $S$. Коэффициенты этих многочленов составляют $j$-й столбец обратной матрицы $A^{-1}$. Действительно, пусть $A^{-1} = (l_{ij})_{i,j=1}^{n+1}$. Тогда
\[
\lambda_j(x) = l_{1j} x_1 + l_{2j} x_2 + \dots + l_{nj} x_n + l_{n+1,j}.
\]

Любой полином $p \in \Pi_1(\mathbb{R}^n)$ может быть разложен по этому базису:
\[
p(x) = \sum_{j=1}^{n+1} p(x^{(j)}) \lambda_j(x). \tag{1}
\]

\subsection{Интерполяционный проектор}

Для заданной функции $f \in C(Q_n)$ определим отображение $P: C(Q_n) \to \Pi_1(\mathbb{R}^n)$ формулой
\[
P f(x) = \sum_{j=1}^{n+1} f(x^{(j)}) \lambda_j(x). \tag{2}
\]

Из свойства $\lambda_j(x^{(k)}) = \delta_{jk}$ следует, что
\[
P f(x^{(j)}) = f(x^{(j)}), \quad j = 1, \dots, n+1,
\]
то есть полином $P f$ интерполирует функцию $f$ в узлах $x^{(j)}$. Отображение $P$ называется \textbf{интерполяционным проектором}.

Оператор $P$ линеен и обладает свойством идемпотентности: для любого полинома $p \in \Pi_1(\mathbb{R}^n)$ выполнено $P p = p$.

\subsection{Норма интерполяционного проектора}

Норма оператора $P$ определяется стандартным образом:
\[
\|P\| = \sup_{\|f\|_{C(Q_n)} \le 1} \|P f\|_{C(Q_n)}.
\]

Из линейности $P$ и того, что образ $P$ конечномерен, можно получить явную формулу для нормы. Поскольку $\lambda_j(x)$ — аффинные функции, а максимум суммы модулей на кубе достигается в одной из вершин, справедливо равенство
\[
\|P\| = \max_{x \in \operatorname{ver}(Q_n)} \sum_{j=1}^{n+1} |\lambda_j(x)|. \tag{3}
\]
% todo: сформулировать
Показано в \cite{nevsky_geo}

\subsection{Неравенство Лебега}

% todo: ссылку на литературу
Пусть $E_n(f)$ — наилучшее приближение функции $f$ полиномами из $\Pi_1(\mathbb{R}^n)$:
\[
E_n(f) = \min_{p \in \Pi_1(\mathbb{R}^n)} \|f - p\|_{C(Q_n)}.
\]

Для интерполяционного проектора $P$ справедливо \textbf{неравенство Лебега}:
\[
\|f - P f\|_{C(Q_n)} \le (1 + \|P\|) \cdot E_n(f). \tag{4}
\]

Это неравенство показывает, что погрешность интерполяции не превосходит наилучшего приближения, умноженного на множитель $(1 + \|P\|)$. Таким образом, чем меньше норма проектора, тем ближе интерполяционный полином к наилучшему приближению.

\subsection{Минимальная норма проектора}

Обозначим через $\theta_n$ минимальную норму интерполяционного проектора при условии, что узлы интерполяции принадлежат кубу $Q_n$ и образуют невырожденный симплекс:
\[
\theta_n = \min_{\substack{x^{(j)} \in Q_n \\ \det(A) \neq 0}} \|P\|.
\]

% todo: что за (3)
Задача о нахождении $\theta_n$ является классической в теории приближений. Она сводится к поиску такого набора узлов (симплекса $S \subset Q_n$), который минимизирует правую часть (3). Отметим, что выражение под знаком максимума зависит от $n(n+1)$ переменных (координат вершин), а построение функций $\lambda_j(x)$ требует обращения матрицы $A$, что делает задачу весьма сложной для больших $n$.

\subsection{Оценки нормы проектора}

% todo: сформулировать
В \cite{nevsky_ukhalov_sel_problems} приведены оценки.... В этой статье приведена таблица ~\ref{tab:theta_estimates} лучших известных оценок для 1 ≤ n ≤ 26.

Для величины $\theta_n$ известны следующие общие оценки (см. \cite{nevsky_geo}):
\[
\frac{1}{4}\sqrt{n} < \theta_n < 3\sqrt{n}, \qquad 3 - \frac{4}{n+1} \le \theta_n.
\]

Если $S$ — симплекс максимального объёма в $Q_n$, а $P$ — интерполяционный проектор с узлами в вершинах $S$, то справедливо асимптотическое равенство
\[
\|P\| \asymp \theta_n,
\]
где запись $f(n) \asymp g(n)$ означает, что существуют константы $c_1, c_2 > 0$, не зависящие от $n$, такие что
\[
c_1 g(n) \le f(n) \le c_2 g(n) \quad \text{для всех достаточно больших } n.
\]

Таким образом, норма проектора по узлам в вершинах симплекса максимального объёма и величина $\theta_n$ имеют одинаковый порядок роста. Это свойство лежит в основе метода уточнения оценок для $\theta_n$, применяемого в настоящей работе. В качестве верхней границы $\theta_n$ можно взять значение нормы проектора с узлами в вершинах симплекса максимального объёма в $Q_n$.
% ============================================================================
\section{Матрицы Адамара}

\subsection{Определение и основные свойства}

\textbf{($-1/1$)-матрицей} называется матрица, каждый элемент которой равен $-1$ или $1$. \textbf{Матрицей Адамара} порядка $n$ называется квадратная $(-1/1)$-матрица $H$ размера $n \times n$, удовлетворяющая условию:
\[
H H^T = n I_n,
\]
где $I_n$ — единичная матрица. Из этого свойства непосредственно следует, что строки (и столбцы) матрицы Адамара попарно ортогональны. Для определителя матрицы Адамара справедливо соотношение:
\[
|\det H| = n^{n/2}.
\]

Подробная информация о матрицах Адамара содержится в книге \cite{hall_comb_1970}

Матрицы Адамара существуют не для всех порядков. Необходимым условием существования является $n = 1$, $n = 2$ или $n$, кратное $4$. Достаточность этого условия составляет содержание \textbf{гипотезы Адамара}, которая до настоящего времени не доказана. Тем не менее, для огромного количества порядков, кратных $4$, матрицы Адамара построены. В частности, для порядка $32 = 4 \cdot 8$ матрицы Адамара существуют.

\subsection{Количество неэквивалентных матриц Адамара}

Две матрицы Адамара называются \textbf{эквивалентными}, если одна может быть получена из другой с помощью перестановок строк и столбцов и умножения строк и столбцов на $-1$. Количество неэквивалентных матриц Адамара для малых порядков приведено в таблице~\ref{tab:hadamard_counts}.

% todo: сформулировать
В таблице .... Из таблицы видно, что количество матриц быстро растёт с увеличением порядка. Для $n = 32$ оно составляет уже более тринадцати миллионов, что делает полный перебор всех симплексов максимального объёма в $\mathbb{R}^{31}$ практически невозможным.

\subsection{Методы построения матриц Адамара}

Существует несколько различных методов построения матриц Адамара. В настоящей работе используются матрицы порядка $32$, полученные следующими способами.

\subsubsection{Метод Сильвестра (тип `syl`)}

Для порядков, являющихся степенью двойки, матрицы Адамара могут быть построены рекурсивным методом Сильвестра. Рекуррентное соотношение имеет вид:
\[
H_{1} = (1), \qquad H_{2n} = \begin{pmatrix}
H_n & H_n \\
H_n & -H_n
\end{pmatrix}.
\]

Например, для $n = 2$ получаем:
\[
H_2 = \begin{pmatrix}
1 & 1 \\
1 & -1
\end{pmatrix},
\]
для $n = 4$:
\[
H_4 = \begin{pmatrix}
1 & 1 & 1 & 1 \\
1 & -1 & 1 & -1 \\
1 & 1 & -1 & -1 \\
1 & -1 & -1 & 1
\end{pmatrix}.
\]

Матрицы, построенные этим методом, обозначаются суффиксом \texttt{syl} (например, \texttt{had.32.syl}).

\subsubsection{Метод Пэли (тип `pal`)}

Конструкция Пэли (Paley construction) позволяет строить матрицы Адамара для порядков $n = q + 1$, где $q$ — степень нечётного простого числа. В основе метода лежат свойства квадратичных вычетов в конечных полях. Для порядка $32 = 31 + 1$ существует матрица Адамара типа Пэли, обозначаемая суффиксом \texttt{pal} (например, \texttt{had.32.pal}).

\subsubsection{Другие типы матриц порядка 32}

Для порядка $32$ кроме матриц типов \texttt{syl} и \texttt{pal} существуют также матрицы, обозначаемые \texttt{t1}, \texttt{t2}, \texttt{t3}, \texttt{t4} (например, \texttt{had.32.t1}, \texttt{had.32.t2}, \texttt{had.32.t3}, \texttt{had.32.t4}). Эти матрицы получены с использованием других конструкций (например, методом Уильямсона или методом Турина). В настоящей работе исследуются все перечисленные типы матриц для сравнения значений норм интерполяционных проекторов, получаемых на соответствующих симплексах.

\subsection{От $(-1/1)$-матриц к $(0/1)$-матрицам максимального определителя}

Как показано в \cite{nevsky_geo}, симплекс максимального объёма в гиперкубе $Q_n = [0,1]^n$ может быть получен из максимального $(0/1)$-определителя порядка $n$. В качестве первых $n$ вершин такого симплекса можно взять строки этого определителя, а в качестве последней $(n+1)$-й вершины добавить нулевую строку.

Максимальный $(0/1)$-определитель порядка $n$ может быть получен из максимального $(-1/1)$-определителя порядка $n+1$ с помощью следующей процедуры. Пусть $A$ — максимальный $(-1/1)$-определитель порядка $n+1$.

\begin{enumerate}
    \item Каждый столбец $A$, начинающийся с $-1$, умножается на $-1$.
    \item Каждая строка новой матрицы, начинающаяся с $-1$, также умножается на $-1$.
\end{enumerate}

Матрицу порядка $n+1$, которая получится в результате выполнения шагов 1–2, обозначим через $B$. У этой $(-1/1)$-матрицы первый столбец и первая строка целиком состоят из $1$.

\begin{enumerate}
    \setcounter{enumi}{2}
    \item Пусть $C$ — подматрица порядка $n$ матрицы $B$, стоящая в строках и столбцах с номерами $2, \dots, n+1$. В этой матрице производится следующая замена элементов: $1$ заменяется на $0$, а $-1$ на $1$. Получившуюся $(0/1)$-матрицу порядка $n$ обозначим $D$.
\end{enumerate}

Полученный таким образом определитель $D$ является максимальным $(0/1)$-определителем порядка $n$.

\subsection{Построение симплекса максимального объёма в $\mathbb{R}^{31}$}

В настоящей работе для размерности $n = 31$ в качестве исходных берутся матрицы Адамара порядка $32$ следующих типов: \texttt{syl}, \texttt{pal}, \texttt{t1}, \texttt{t2}, \texttt{t3}, \texttt{t4}. Для каждой из них выполняется описанная выше процедура перехода к $(0/1)$-матрице $D$ порядка $31$, после чего строки $D$ становятся первыми $31$ вершиной симплекса, а в качестве последней, $32$-й вершины добавляется нулевая строка:
\[
V = \{0, \operatorname{row}_1(D), \operatorname{row}_2(D), \dots, \operatorname{row}_{31}(D)\} \subset [0,1]^{31}.
\]

Таким образом, для каждого типа матрицы Адамара получается соответствующий симплекс максимального объёма, вписанный в гиперкуб $[0,1]^{31}$. Все вершины таких симплексов являются вершинами гиперкуба. Полученные симплексы используются в настоящей работе для вычисления верхних оценок минимальной нормы интерполяционного проектора $\theta_{31}$.

\subsection{Замечание о количестве симплексов}

Для размерности $n = 31$ количество неэквивалентных матриц Адамара порядка $32$ составляет $13\,710\,027$. Это означает, что существует такое же количество различных симплексов максимального объёма, вписанных в гиперкуб $Q_{31}$. Значения норм интерполяционных проекторов для разных симплексов могут различаться.

% todo: сформулировать
В таблице  ~\ref{tab:hadamard_counts} приводятся данные с \cite{koukouvinos_hadamard}....

В настоящей работе для исследования отобраны шесть матриц Адамара порядка $32$, доступных на сайте библиотеки Нила Слоуна \cite{sloane_hadamard}:
\[
\texttt{had.32.syl},\quad \texttt{had.32.pal},\quad \texttt{had.32.t1},\quad \texttt{had.32.t2},\quad \texttt{had.32.t3},\quad \texttt{had.32.t4}.
\]

Матрица типа \texttt{syl} построена рекурсивным методом Сильвестра, матрица типа \texttt{pal} получена с помощью конструкции Пэли, а матрицы типов \texttt{t1} – \texttt{t4} представляют собой другие известные неэквивалентные матрицы Адамара порядка $32$.

Для каждой из этих шести матриц строится соответствующий симплекс максимального объёма по описанной выше процедуре, после чего вычисляется норма интерполяционного проектора. Полученные значения дают верхние оценки для $\theta_{31}$.

Полный перебор всех $13\,710\,027$ симплексов является крайне трудоёмкой задачей и не входит в рамки настоящей работы. Исследование ограничивается указанными шестью представителями, что позволяет получить содержательные оценки при разумных вычислительных затратах.
% ============================================================================
\section{Описание программной реализации}

Программная реализация выполнена на языке C++ (стандарт C++17). Для линейно-алгебраических операций используется библиотека Eigen \cite{eigen_lib}. Сборка проекта осуществляется с помощью CMake.

\subsection{Общая архитектура}

Программа построена в соответствии с принципами SOLID. Основные классы:

\begin{itemize}
    \item \texttt{ISimplexProvider} — интерфейс для получения вершин симплекса;
    \item \texttt{SilvesterMethodProvider} — построение симплекса из матрицы Адамара Сильвестра;
    \item \texttt{HadamardFileProvider} — чтение матриц Адамара из файлов;
    \item \texttt{SimplexProviderFactory} — фабрика для создания провайдеров;
    \item \texttt{HadamardMatrixIterator} — итератор для последовательного чтения матриц из файла;
    \item \texttt{LebesgueFunction} — вычисление константы Лебега в точке;
    \item \texttt{FullEnumUpperBoundCalculator} — многопоточный перебор вершин гиперкуба;
    \item \texttt{CommandLineParser} — разбор аргументов командной строки;
    \item \texttt{Processing} — организация основных вычислительных режимов.
\end{itemize}

\subsection{Ввод и вывод данных}

Программа поддерживает два режима работы:
\begin{enumerate}
    \item Построение симплекса по матрице Адамара Сильвестра (по умолчанию);
    \item Загрузка матриц Адамара из файла (символьный или числовой формат);
\end{enumerate}

Файлы с матрицами Адамара могут быть загружены с сайта библиотеки Нила Слоуна \cite{sloane_hadamard}. Для автоматической загрузки всех доступных матриц порядка 32 разработан скрипт \texttt{download\_hadamard.sh}. Он скачивает файлы \texttt{had.32.syl}, \texttt{had.32.pal}, \texttt{had.32.t1} – \texttt{had.32.t4} и объединяет их в один файл с пустыми строками-разделителями.

Управление режимами работы осуществляется через аргументы командной строки:
\begin{itemize}
    \item \texttt{--dimension N} — установить размерность пространства;
    \item \texttt{--load-hadamard FILE} — загрузить матрицы Адамара из файла;
    \item \texttt{--batch DIR} — пакетная обработка всех файлов в каталоге;
    \item \texttt{--threads N} — количество потоков (0 = автоматически);
    \item \texttt{--log-interval N} — интервал логирования (количество вершин);
    \item \texttt{--range START END} — частичный перебор вершин в заданном диапазоне.
\end{itemize}


\subsection{Дополнительные проверки корректности}

Для верификации программной реализации и достоверности полученных результатов был проведён ряд дополнительных проверок.

\subsubsection{Проверка интерполяционных свойств}

Для каждого построенного симплекса выполнялась проверка условия
\[
\lambda_j(x^{(k)}) = \delta_{jk}, \qquad j,k = 1,\dots,n+1,
\]
где $\lambda_j$ — базисные многочлены Лагранжа (барицентрические координаты), а $x^{(k)}$ — вершины симплекса. Проверка проводилась с точностью $10^{-5}$. Для всех рассмотренных матриц Адамара максимальная ошибка не превысила $10^{-10}$, что подтверждает корректность построения симплексов и вычисления базисных функций.

\subsubsection{Проверка определителей матриц Адамара}

Для каждой загруженной матрицы Адамара $H$ порядка $n$ вычислялся определитель и сравнивался с теоретическим значением $n^{n/2}$. Относительная погрешность не превышала $10^{-10}$, что подтверждает, что загруженные матрицы действительно являются матрицами Адамара.

\subsubsection{Проверка полного перебора}

Для малых размерностей ($n = 3, 7, 11, 15, 19, 23, 27$) был выполнен полный перебор всех вершин гиперкуба, и полученные значения константы Лебега были сравнены с известными из литературы. Результаты сравнения соответствуют приведенным в ~\ref{tab:theta_estimates}.

\subsubsection{Повторная проверка найденного максимума}

После завершения полного перебора для каждой матрицы значение константы Лебега повторно вычислялось в найденной вершине. Расхождение между значением, зафиксированным во время перебора, и значением, полученным при повторной проверке, не превышало $10^{-10}$, что подтверждает корректность синхронизации глобального максимума при многопоточном переборе.

\subsubsection{Сравнение с известными результатами для малых размерностей}

Для размерности $n = 7$ был построен симплекс из матрицы Адамара Сильвестра порядка $8$, и вычисленная константа Лебега составила $\|P\| = 2.82843$, что совпадает с теоретической оценкой $\sqrt{8}$. Это дополнительно подтверждает корректность программной реализации.

\subsection{Вычисление нормы проектора}

Норма интерполяционного проектора вычисляется по формуле (2.3):
\[
\|P\| = \max_{x \in \operatorname{ver}(Q_n)} \sum_{j=1}^{n+1} |\lambda_j(x)|.
\]

Для нахождения максимума реализован полный перебор всех $2^{31}$ вершин гиперкуба. Каждая вершина кодируется 31-битовым целым числом, что позволяет эффективно генерировать координаты с помощью битовых операций.

Вычисления распараллелены на несколько потоков (количество задаётся параметром \texttt{--threads}). Для синхронизации глобального максимума используются атомарные операции. Программа выводит прогресс выполнения каждые $10^7$ вершин с оценкой времени до завершения.

По умолчанию все результаты выводятся в стандартный поток вывода (консоль). Для сохранения результатов в файл при одновременном отображении в консоли может быть использована утилита \texttt{tee}, например:
\begin{verbatim}
./HadamardSimplex --batch preset/ | tee results.txt
\end{verbatim}
При указании опции \texttt{--output FILE} программа самостоятельно сохраняет результаты в указанный файл.

Полный исходный код программы приведён в приложении B.
% ============================================================================
\section{Результаты вычисления нормы}

\subsection{Параметры вычислительного эксперимента}

Вычисления проводились на следующем оборудовании и программном обеспечении:

\begin{itemize}
    \item \textbf{Процессор}: 12th Gen Intel(R) Core(TM) i5-1240P (16 ядер)
    \item \textbf{Оперативная память}: 15 ГБ
    \item \textbf{Операционная система}: Ubuntu 24.04.3 LTS
    \item \textbf{Компилятор}: GCC 13.3.0
    \item \textbf{Библиотека}: Eigen 3
    \item \textbf{Флаги компиляции}: \texttt{-O3 -march=native}
    \item \textbf{Количество потоков}: 16
    \item \textbf{Интервал логирования}: $10^7$ вершин
\end{itemize}

Полный перебор всех $2^{31}$ вершин гиперкуба занял 290 секунд (около 5 минут). Средняя скорость обработки составила около 7.4 миллионов вершин в секунду.

\subsection{Результаты для размерности $n = 31$}

В результате вычислительного эксперимента для симплекса, построенного на основе матрицы Адамара Сильвестра, получено следующее значение максимальной константы Лебега:
\[
\|P\| = 5.000000.
\]

Соответственно, для минимальной нормы интерполяционного проектора $\theta_{31}$ получена верхняя оценка:
\[
\theta_{31} \le 5.000000.
\]

\subsection{Сравнение с теоретической оценкой}

Теоретическая верхняя граница для минимальной нормы интерполяционного проектора в размерности $n = 31$ составляет:
\[
\theta_{31} \le \sqrt{32} \approx 5.656854249.
\]

Полученное в работе значение $\theta_{31} \le 5.000000$ ниже теоретической оценки. Это подтверждает корректность разработанного алгоритма и демонстрирует, что для выбранной матрицы Адамара (типа \texttt{syl}) норма проектора оказалась меньше теоретической границы.

\subsection{Сравнение с результатами для других типов матриц}

В работе были также исследованы матрицы Адамара других типов: \texttt{pal}, \texttt{t1}, \texttt{t2}, \texttt{t3}, \texttt{t4}. Для каждой из них был построен соответствующий симплекс и вычислена максимальная константа Лебега. Результаты приведены в таблице~\ref{tab:hadamard_types_results}.

\begin{table}[h]
\centering
\caption{Значения максимальной константы Лебега для различных типов матриц Адамара порядка 32}
\label{tab:hadamard_types_results}
\begin{tabular}{|c|c|}
\hline
\textbf{Тип матрицы} & $\|P\|$ \\
\hline
\texttt{syl} & 5.000000 \\
\texttt{pal} & $[\text{значение}]$ \\
\texttt{t1} & $[\text{значение}]$ \\
\texttt{t2} & $[\text{значение}]$ \\
\texttt{t3} & $[\text{значение}]$ \\
\texttt{t4} & $[\text{значение}]$ \\
\hline
\end{tabular}
\end{table}

Как видно из таблицы, значения норм проекторов для разных типов матриц могут различаться. Наименьшее значение среди рассмотренных типов составило $[\text{значение}]$.

\subsection{Обсуждение результатов}

Полученная оценка $\theta_{31} \le 5.000000$ является новой и уточняет ранее известные теоретические границы. Для размерности $n = 31$ ранее не было известно численных значений верхней границы минимальной нормы интерполяционного проектора. Результат, полученный в настоящей работе, может быть использован для дальнейшего уточнения оценок $\theta_n$ и для верификации других численных методов.

Следует отметить, что исследование ограничилось шестью типами матриц Адамара. Полный перебор всех $13\,710\,027$ неэквивалентных матриц порядка $32$ является крайне трудоёмкой задачей и не входил в рамки настоящей работы. Возможно, среди оставшихся матриц существуют симплексы, дающие ещё меньшие значения нормы проектора, что может быть предметом дальнейших исследований.
% ============================================================================
\section*{Заключение}
\addcontentsline{toc}{section}{Заключение}

В ходе выполнения выпускной квалификационной работы была разработана программа на языке C++, предназначенная для вычисления верхних оценок минимальной нормы интерполяционного проектора $\theta_n$ для симплексов, построенных из матриц Адамара. Программа поддерживает:

\begin{itemize}
    \item построение симплекса из матрицы Адамара Сильвестра;
    \item загрузку матриц Адамара из файлов (символьный и числовой форматы);
    \item многопоточный полный перебор всех $2^{31}$ вершин гиперкуба;
    \item логирование прогресса с оценкой времени до завершения;
    \item сохранение результатов в выходной файл.
\end{itemize}

\subsection*{Полученные результаты}

В результате вычислительного эксперимента для симплекса, построенного на основе матрицы Адамара Сильвестра порядка $32$, получена верхняя оценка
\[
\theta_{31} \le 5.000000.
\]
Теоретическая верхняя граница для данной размерности составляет $\sqrt{32} \approx 5.656854249$, что подтверждает корректность разработанного алгоритма.

\subsection*{Трудоёмкость анализа следующей размерности}

Следующая размерность, для которой существует матрица Адамара, — $n = 63$ (матрица порядка $64$). В этом случае гиперкуб содержит $2^{63} \approx 9.22 \times 10^{18}$ вершин. Даже при сохранении скорости обработки $7.4 \times 10^6$ вершин в секунду полный перебор занял бы более $80\,000$ лет, что делает его практически невозможным. Таким образом, разработанный метод полного перебора применим только для размерностей, не превышающих $31$. Для анализа более высоких размерностей требуется разработка новых подходов, например, случайной выборки вершин или использования теоретических оценок.

\subsection*{Перспективы дальнейших исследований}

В работе были исследованы шесть типов матриц Адамара порядка $32$ (\texttt{syl}, \texttt{pal}, \texttt{t1}, \texttt{t2}, \texttt{t3}, \texttt{t4}). Полный перебор всех $13\,710\,027$ неэквивалентных матриц является крайне трудоёмкой задачей и не входил в рамки настоящей работы. Перспективным направлением дальнейших исследований является:

\begin{itemize}
    \item анализ других типов матриц Адамара порядка $32$;
    \item разработка эвристических методов поиска симплексов с малой нормой проектора без полного перебора вершин;
    \item применение статистических методов (случайная выборка вершин) для оценки $\theta_n$ при $n > 31$;
    \item исследование возможности использования квантовых алгоритмов (например, алгоритма Гровера) для ускорения перебора.
\end{itemize}

\subsection*{Заключительное замечание}

Разработанная программа и полученные результаты могут быть использованы для дальнейшего уточнения оценок $\theta_n$, а также для верификации других численных методов в теории интерполяционных проекторов. Подготовленный набор матриц Адамара порядка $32$ в машинно-читаемой форме опубликован в открытом репозитории.

% ============================================================================

\begin{table}[h]
\centering
\caption{Количество неэквивалентных матриц Адамара}
\label{tab:hadamard_counts}
\begin{tabular}{|c|c|}
\hline
\textbf{Порядок $n$} & \textbf{Количество неэквивалентных матриц} \\
\hline
1 & 1 \\
2 & 1 \\
4 & 1 \\
8 & 1 \\
12 & 1 \\
16 & 5 \\
20 & 3 \\
24 & 60 \\
28 & 487 \\
32 & 13\,710\,027 \\
\hline
\end{tabular}
\end{table}

\begin{table}[h]
\centering
\caption{Верхние оценки $\theta_n$ для $1 \le n \le 27$}
\label{tab:theta_estimates}
\begin{tabular}{|c|c|c|}
\hline
$n$ & $\theta_n \le$ & \textbf{Источник} \\
\hline
1 & 1 & AH \\
2 & $\sqrt{5} \approx 2.2361$ & A \\
3 & $2.5$ & A \\
4 & $\frac{13}{5} = 2.6$ & A \\
5 & $2.5$ & A \\
6 & $3.1428$ & M \\
7 & $2.5$ & AH \\
8 & $3.1428$ & M \\
9 & $3.8$ & M \\
10 & $3.8$ & M \\
11 & $3.7692$ & H \\
12 & $3.4$ & M \\
13 & $3.7692$ & H \\
14 & $4.2$ & M \\
15 & $4.2$ & M \\
16 & $4.0882$ & H \\
17 & $5.5882$ & M \\
18 & $4.2$ & M \\
19 & $4.7241$ & H \\
20 & $4.7241$ & H \\
21 & $5.02$ & M \\
22 & $5.4238$ & M \\
23 & $4.9047$ & M \\
24 & $4.9047$ & M \\
25 & $5.2087$ & M \\
26 & $5.2087$ & M \\
27 & $5.0$ & UH \\
\hline
\multicolumn{3}{l}{\footnotesize\textit{Примечание:} A — аналитическая оценка, H — оценка, полученная с помощью компьютера} \\
\multicolumn{3}{l}{\footnotesize\textit{(Хлесткова, Невский), M — оценка из работы \cite{nevsky_ukhalov_2015},}} \\
\multicolumn{3}{l}{\footnotesize\textit{UH — оценка, полученная в настоящей работе.}} \\
\end{tabular}
\end{table}

\printbibliography

% ============================================================================
\appendix
\section{Руководство пользователя программы}
\label{app:user_guide}

% Описание параметров командной строки, форматов входных файлов.

\hfill [Заполнить]

\section{Листинги основных модулей программы}
\label{app:listings}

% Код основных классов: HadamardMatrixIterator, FullEnumUpperBoundCalculator и др.

\hfill [Заполнить]

\section{Дополнительные материалы}
\label{app:additional}

% Другие материалы по необходимости.

\hfill [Заполнить]

\end{document}

\end{document}